\section{Evaluation Results}
\begin{table}[t]
\centering
\scriptsize
\begin{tabular}{@{}lrrrr@{}}
\toprule
Retriever & Hit@1 & Hit@5 & MRR & nDCG@5 \\
\midrule
Keyword & .71 & .91 & .796 & .757 \\
BM25 & .84 & .97 & .892 & .855 \\
Dense text & .76 & .95 & .838 & .796 \\
BM25+dense RRF & \textbf{.85} & \textbf{.98} & \textbf{.905} & .853 \\
RRF+modality & .59 & .95 & .734 & .714 \\
\bottomrule
\end{tabular}
\caption{\textbf{The modality heuristic harms first-page retrieval.} Retrospective page-level results on 100 questions/25 papers. Hit@k means any gold page at k, not fractional recall.}
\label{tab:retrieval}
\end{table}

BM25+dense improves Hit@1 over BM25 by only .01 (paired 95\% CI $[-.054,.078]$), so its incremental benefit is unresolved. Adding the modality heuristic lowers Hit@1 by .26 (CI $[-.388,-.136]$) and MRR by .172 (CI $[-.253,-.092]$) relative to BM25+dense. This negative ablation is consequential: extra visual intent and structural boosts can displace a correct first page even while Hit@5 remains relatively high. It does not isolate which sub-rule caused the loss.

\begin{table}[t]
\centering
\scriptsize
\begin{tabular}{@{}lrrr@{}}
\toprule
Automatic diagnostic & Off & On & $\Delta$ (95\% CI) \\
\midrule
Token F1 & .047 & .067 & +.020 [.015,.027] \\
Ref.-unit F1 & .159 & .224 & +.065 [.042,.091] \\
Key-term coverage & .759 & .742 & $-$.017 [$-$.038,.000] \\
Cited-page alignment & .670 & .618 & $-$.052 [$-$.091,$-$.015] \\
\bottomrule
\end{tabular}
\caption{\textbf{Repair improves lexical overlap but worsens measured page alignment.} Thirty paired answerable cases; intervals are paper-clustered paired bootstraps. Page alignment is defined for only the 20 cases with a gold page label. All metrics are automatic proxies.}
\label{tab:repair}
\end{table}

Repair changes the final text in all 30 answerable cases. Its overlap gains are small in absolute terms and do not establish human-judged factual support. The 10 remaining cases name a figure or table as the gold locus but provide no gold page, so treating their unscored alignment as zero would be invalid. The pipeline records 28/30 \textsc{Success} without repair and 30/30 with repair, but success is not correctness.

The missing \emph{gold page labels} are concentrated entirely in the 10 visual questions: none has a paired page-alignment value, whereas all 10 mathematical and all 10 textual questions do. All 10 visual outputs have page-bearing citations in both conditions; those pages cannot be checked against a gold visual locus with this annotation. Exploratory category-level overlap gains appear in all three groups (Appendix Table~\ref{app:subgroups}), but they do not establish visual support.

For the 20 cross-paper controls, \textsc{Abstained} falls from 13/20 without repair to 5/20 with repair (paired change $-$.40, CI $[-.619,-.188]$). Four cases in each condition abstain before generation. Some \textsc{Success} outputs nevertheless refuse in prose, while others discuss unrelated target-paper content; status is thus not a valid semantic-refusal measure. The observed direction is adverse for a status-based safety story, and no hallucination-reduction claim follows.

\begin{table}[t]
\centering
\scriptsize
\begin{tabular}{@{}p{.24\columnwidth}p{.68\columnwidth}@{}}
\toprule
LLM-judge aspect & Recorded distribution \\
\midrule
Correctness ($n=113$) & 59 correct; 33 partly correct; 20 incorrect; 1 not judgeable \\
Grounding ($n=150$) & 112 supported; 26 partly supported; 10 unsupported; 2 not judgeable \\
Citation quality ($n=150$) & 110 accurate; 28 partly accurate; 10 inaccurate; 2 missing \\
\bottomrule
\end{tabular}
\caption{\textbf{The LLM judge identifies both answer and evidence failures.} Correctness uses only reference-eligible cases; 37 additional cases failed the reference gate. Counts are post-hoc model labels, not human ratings.}
\label{tab:llm}
\end{table}

Table~\ref{tab:llm} reports the completed LLM-based evaluation without turning its ordinal labels into a single score. One reference-eligible case was still marked not judgeable because the recorded reason incorrectly called its reference ineligible; we preserve that outcome rather than recoding it. Across all 150 cases, the most frequent non-\texttt{none} failure label is \texttt{missing\_fact} (53 cases), followed by \texttt{wrong\_fact} (14). These labels are mutually exclusive judgments from one LLM, not error rates with independent adjudication. The visual labels are omitted for the protocol mismatch described in Section~5.

\paragraph{Human evaluation analysis (pending).}
At the September 21, 2026 manuscript freeze, the reserved 50-question study has 0 of 250 planned ratings and 0 of 50 cases with any rating. Accordingly, there are no human means, confidence intervals, agreement estimates, reference-verification outcomes, or evidence-support counts to analyze. This paragraph is the explicit insertion point for those planned outputs after all five independently completed files have been validated and imported; until then, Table~\ref{tab:llm} must not be described as human evaluation.

Table~\ref{tab:latency} profiles 100 valid requests; end-to-end p90 is 41.61\,s. Generation dominates the median, while first-use index construction contributes to the retrieval tail. We observed no terminal errors in the corrected batches. We did not measure peak process memory or controlled cold/warm latency; a batch-order difference cannot be interpreted as a repair speedup.

\begin{table}[t]
\centering
\small
\begin{tabular}{lrr}
\toprule
Measured duration & p50 (s) & p95 (s) \\
\midrule
End-to-end & 27.34 & 149.54 \\
Retrieval stage & 0.52 & 112.85 \\
Generation stage & 26.59 & 37.48 \\
\bottomrule
\end{tabular}
\caption{\textbf{Generation dominates the median; first-use retrieval contributes to the tail.} One local workstation, 100 valid requests, uncontrolled cache state. Stage quantiles need not add to end-to-end quantiles.}
\label{tab:latency}
\end{table}

Taken together, the results distinguish what is inspectable from what is supported: the trace exposes each decision, while the negative ablations and incomplete human evidence constrain the conclusions. The next section turns those constraints into deployment lessons.
